\documentclass[11pt]{amsart}

\usepackage[T1]{fontenc}
\usepackage{lmodern}
\usepackage{microtype}
\usepackage{amsmath,amssymb,amsthm}
\usepackage{booktabs}
\usepackage[hidelinks]{hyperref}

\hypersetup{
  pdftitle={Vu's Cover-Ideal Linearity Question Through Twelve Nonisolated Vertices},
  pdfsubject={A computer-assisted finite verification for powers of bipartite cover ideals},
  pdfkeywords={v-number, cover ideal, bipartite graph, exhaustive verification}
}

\newtheorem{theorem}{Theorem}
\newtheorem{proposition}[theorem]{Proposition}
\newtheorem{lemma}[theorem]{Lemma}

\DeclareMathOperator{\Ass}{Ass}
\newcommand{\M}{\mathcal M}

\title[Vu's Linearity Question Through Twelve Vertices]
{Vu's Cover-Ideal Linearity Question Through Twelve Nonisolated Vertices}
\date{}

\subjclass[2020]{Primary 13F55; Secondary 05C69}
\keywords{$v$-number, cover ideal, bipartite graph, exhaustive verification}

\begin{document}

\begin{abstract}
Vu asked whether $v(J(G)^t)$ is linear in $t$ for every bipartite graph $G$.
We prove by exact exhaustive computation that the answer is affirmative for every
bipartite graph with at least one edge and at most twelve nonisolated vertices:
\[
  v\bigl(J(G)^t\bigr)=v(J(G))+(t-1)\tau(G),\qquad t\geq1.
\]
Vu's local formula reduces the problem to an edge condition involving maximum
independent sets. We verify the stronger statement that every edge maximizing
$\alpha(G-N[e])$ has an endpoint in a maximum independent set. The computation
checks $465{,}766{,}775$ neighborhood-multiset encodings using exact integer
arithmetic. This is a finite verification; the general question remains open.
\end{abstract}

\maketitle

\section{Result}

Let $G$ be a finite simple graph, let $S=\Bbbk[x_v:v\in V(G)]$, and let
\[
  J(G)=\bigcap_{uv\in E(G)}(x_u,x_v)
\]
be its cover ideal. For a nonzero homogeneous ideal $I\subsetneq S$ and
$P\in\Ass(S/I)$, set
\begin{align*}
  v_P(I)&=\min\{\deg f:(I:f)=P,\ f\in S\text{ homogeneous}\},\\
  v(I)&=\min_{P\in\Ass(S/I)}v_P(I).
\end{align*}
Write $\tau(G)$ and $\alpha(G)$ for the vertex-cover and independence numbers.

\begin{theorem}\label{thm:main}
Let $G$ be a finite simple bipartite graph with at least one edge and at most
twelve nonisolated vertices. Then, for every integer $t\geq1$,
\[
  v\bigl(J(G)^t\bigr)=v(J(G))+(t-1)\tau(G).
\]
\end{theorem}

Vu proved local linearity \cite[Theorem~1.3]{Vu} and global linearity for
$t\geq |V(G)|-2$ \cite[Lemma~3.1]{Vu}, and asked whether global linearity holds
for every $t\geq1$ \cite[Question~3.2]{Vu}. That question is posed immediately
after the report that no bipartite graph with a nonlinear $v$-function of the
powers of its cover ideal had been found \cite[remark preceding
Question~3.2]{Vu}, so the contribution of the present note is to make that
search exhaustive through twelve nonisolated vertices rather than to introduce
the first such search. Hang and Vu proved the tree case
\cite[Theorem~1.2]{HangVu}. Theorem \ref{thm:main} is a finite verification of
Vu's question, not a resolution for arbitrary bipartite graphs.

\section{Reduction to an edge condition}

We first dispose of isolated vertices.

\begin{lemma}\label{lem:variables}
Let $R$ be a standard graded polynomial ring, let $I\subseteq R$ be homogeneous,
and put $T=R[y_1,\ldots,y_r]$. Then
\[
  \Ass_T(T/IT)=\{QT:Q\in\Ass_R(R/I)\},
\]
and $v_{PT}(IT)=v_P(I)$ for every $P\in\Ass_R(R/I)$.
\end{lemma}

\begin{proof}
The associated-prime identity follows from flat base change. The inequality
$v_{PT}(IT)\leq v_P(I)$ is immediate. Conversely, let
$F=\sum_{\mathbf a}f_{\mathbf a}y^{\mathbf a}$ be a homogeneous witness for
$PT$. Coefficient comparison gives
\[
  P=(IT:F)\cap R=\bigcap_{f_{\mathbf a}\neq0}(I:f_{\mathbf a}),
  \qquad P\subseteq(I:f_{\mathbf a}).
\]
Modulo $P$, a finite intersection of nonzero ideals in the domain $R/P$ is
nonzero. Hence $(I:f_{\mathbf a})=P$ for some coefficient, and
$\deg f_{\mathbf a}\leq\deg F$.
\end{proof}

We may therefore assume that $G$ has no isolated vertices. For a vertex $u$,
let $N[u]$ be its closed neighborhood; for an edge $e=uv$, put
$N[e]=N[u]\cup N[v]$ and $P_e=(x_u,x_v)$. Since a graph is unimodular exactly
when it is bipartite, standard gradedness of the vertex-cover algebra gives
$J(G)^t=J(G)^{(t)}$ for every $t\geq1$ \cite[Theorem~1.1]{HHT}. Thus
\[
  J(G)^t=\bigcap_{e\in E(G)}P_e^t.
\]
The edge primes are pairwise incomparable, so each $P_e$ is a minimal prime of
$J(G)^t$ and hence lies in $\Ass(S/J(G)^t)$, while the displayed intersection of
$P_e$-primary ideals gives the reverse inclusion. Thus
$\Ass(S/J(G)^t)=\{P_e:e\in E(G)\}$ and
\[
  v\bigl(J(G)^t\bigr)=\min_{e\in E(G)}v_{P_e}\bigl(J(G)^t\bigr).
\]

For $e=uv$, Vu's formula \cite[Theorem~1.3 and Lemmas~2.5--2.6]{Vu} becomes
\begin{align}
  v_{P_e}\bigl(J(G)^t\bigr)&=c_e+(t-1)g_e, \label{eq:local}\\
  c_e&=\deg(u)+\deg(v)-2+\tau(G-N[e]), \label{eq:c}\\
  g_e&=\min\{h(u),h(v)\},\qquad h(u)=\deg(u)+\tau(G-N[u]). \label{eq:g}
\end{align}
The common-neighbor term in Vu's general expression for $c_e$ vanishes because
$G$ is bipartite.

\begin{lemma}\label{lem:graph}
Let $G$ be bipartite without isolated vertices, let $n=|V(G)|$, and let
$e=uv\in E(G)$. Then
\[
  c_e=n-2-\alpha(G-N[e]).
\]
Moreover, $h(u)\geq\tau(G)$, with equality if and only if $u$ belongs to a
maximum independent set of $G$, and
\[
  \min_{e\in E(G)}g_e=\tau(G).
\]
\end{lemma}

\begin{proof}
Since $|N[e]|=\deg(u)+\deg(v)$ in a bipartite graph, equation \eqref{eq:c} and
$\tau(H)=|V(H)|-\alpha(H)$ give the first identity. The quantity $h(u)$ is the
minimum size of a vertex cover avoiding $u$: such a cover must contain $N(u)$
and then cover $G-N[u]$. Hence $h(u)=\tau(G)$ exactly when a minimum vertex
cover avoids $u$, equivalently when a maximum independent set contains $u$.
Finally, choose a vertex $u$ in a maximum independent set and an incident edge
$uv$. Then $g_{uv}\leq h(u)=\tau(G)$, while every $g_e\geq\tau(G)$.
\end{proof}

Define
\[
  \M(G)=\operatorname*{arg\,max}_{e\in E(G)}\alpha(G-N[e]).
\]

\begin{proposition}[Edge criterion]\label{prop:criterion}
For a bipartite graph $G$ without isolated vertices, the following are
equivalent:
\begin{enumerate}
  \item $t\mapsto v(J(G)^t)$ is affine for every integer $t\geq1$;
  \item $v(J(G)^t)=v(J(G))+(t-1)\tau(G)$ for every integer $t\geq1$;
  \item some edge in $\M(G)$ has an endpoint in a maximum independent set of
  $G$.
\end{enumerate}
\end{proposition}

\begin{proof}
By \eqref{eq:local}, the global $v$-number is the lower envelope of finitely many
lines $c_e+(t-1)g_e$. Its value at $t=1$ is $\min_e c_e=v(J(G))$, and its
eventual slope is $\min_e g_e=\tau(G)$ by Lemma \ref{lem:graph}. A finite lower
envelope is affine on all positive integers if and only if one line minimizes
both its intercept and its slope: in the nontrivial direction, a line attaining
the envelope for infinitely many integers agrees with the affine envelope at
two points and hence identically. Lemma \ref{lem:graph} identifies the
intercept minimizers with $\M(G)$ and the slope minimizers with the edges having
an endpoint in a maximum independent set.
\end{proof}

\section{Exact exhaustive verification}\label{sec:verification}

Fix $n\leq12$, choose $1\leq q\leq\lfloor n/2\rfloor$, and put $p=n-q$.
Every bipartite graph without isolated vertices is represented, possibly more
than once, by a tuple
\[
  (A_1,\ldots,A_q),
  \qquad
  \varnothing\neq A_i\subseteq[p],
  \qquad
  A_1\cup\cdots\cup A_q=[p],
\]
where $A_i$ is the neighborhood of the $i$th vertex on the $q$-side, and the
tuple is nondecreasing with respect to the integer value of the characteristic
bitmask of $A_i$ in $[p]$.

Let $L=[p]$, $R=[q]$, and $U(B)=\bigcup_{i\in B}A_i$ for $B\subseteq R$.
For every encoded graph,
\begin{equation}\label{eq:alpha}
  \alpha(G)=\max_{B\subseteq R}\bigl(|B|+p-|U(B)|\bigr).
\end{equation}
Indeed, once the right part $B$ of an independent set is fixed, every vertex of
$L\setminus U(B)$ may be added. If $e=xi$ with $x\in L$ and $i\in R$, then
\begin{equation}\label{eq:edgealpha}
  \alpha(G-N[e])
  =\max_{B\subseteq R\setminus N(x)}
  \bigl(|B|+p-|A_i\cup U(B)|\bigr).
\end{equation}
The maximizers in \eqref{eq:alpha} also determine membership in maximum
independent sets: $i\in R$ occurs in one if and only if some maximizing $B$
contains $i$, while $x\in L$ occurs in one if and only if some maximizing $B$
satisfies $x\notin U(B)$. Thus \eqref{eq:alpha}--\eqref{eq:edgealpha} give an
exact bitset test of the condition in Proposition \ref{prop:criterion}.

For fixed $(p,q)$, the number of tuples before imposing full union is
\[
  E(p,q)=\binom{2^p+q-2}{q},
\]
and inclusion--exclusion gives the retained count
\[
  K(p,q)=\sum_{r=0}^{p}(-1)^r\binom{p}{r}
  \binom{2^{p-r}+q-2}{q},
\]
with $\binom ab=0$ for $a<b$. The generated counts agree with these formulas
for every $(p,q)$.

\begin{table}[ht]
\centering
\small
\caption{Complete neighborhood-multiset census. In no retained encoding does an
edge maximizing $\alpha(G-N[e])$ fail to have an endpoint in a maximum
independent set; see Proposition \ref{prop:finite} below.}
\label{tab:census}
\begin{tabular}{r@{\qquad}r@{\qquad}r}
\toprule
$n$ & generated & retained \\
\midrule
 2 &           1 &           1 \\
 3 &           3 &           1 \\
 4 &          13 &           5 \\
 5 &          43 &          14 \\
 6 &         235 &          98 \\
 7 &       1{,}239 &         522 \\
 8 &      10{,}659 &       5{,}472 \\
 9 &      98{,}439 &      53{,}733 \\
10 &   1{,}428{,}007 &     895{,}828 \\
11 &  23{,}944{,}527 &  16{,}033{,}665 \\
12 & 610{,}303{,}679 & 448{,}777{,}436 \\
\midrule
Total & 635{,}786{,}845 & 465{,}766{,}775 \\
\bottomrule
\end{tabular}
\end{table}

\begin{proposition}\label{prop:finite}
Let $G$ be a bipartite graph without isolated vertices and with
$|V(G)|\leq12$. Every edge in $\M(G)$ has an endpoint belonging to a maximum
independent set of $G$.
\end{proposition}

\begin{proof}
The enumeration above covers every such graph. The computation evaluates
\eqref{eq:alpha} and \eqref{eq:edgealpha} using integer bit operations and checks
every edge in $\M(G)$. None of the $465{,}766{,}775$ retained encodings violates
the stated condition.
\end{proof}

\begin{proof}[Proof of Theorem \ref{thm:main}]
Delete isolated vertices using Lemma \ref{lem:variables}. Proposition
\ref{prop:finite} supplies an edge satisfying Proposition
\ref{prop:criterion}(3), which yields the asserted formula.
\end{proof}

\paragraph{Exact verification.}
Everything the computation does is fixed by the data printed above: the encoding
of Section \ref{sec:verification}, the closed forms $E(p,q)$ and $K(p,q)$, the
two bitset formulas \eqref{eq:alpha}--\eqref{eq:edgealpha}, and the rule that
reads membership in a maximum independent set off the maximizers of
\eqref{eq:alpha}. All arithmetic is exact integer and bit-mask arithmetic, with
no floating point and no randomness. The computation was re-run once
independently, on a separate machine, by a reimplementation written from that
data alone, and all $53$ of its checks agree with what is reported here: both
columns of Table \ref{tab:census} and both totals come out three ways, from the
closed forms, from an unbounded-knapsack dynamic program and from prefix
enumeration; equations \eqref{eq:alpha}--\eqref{eq:edgealpha}, the membership
rule, Lemma \ref{lem:graph} and Propositions \ref{prop:criterion} and
\ref{prop:finite} agree with brute force on all $59{,}846$ retained encodings
with $n\leq9$; and the condition of Proposition \ref{prop:finite} fails on
$4{,}381$ of the $28{,}262$ labelled graphs without isolated vertices on three
to six vertices once bipartiteness is dropped, while no bipartite graph among
those $28{,}262$ fails it. The
$955{,}674$ retained encodings with $n\leq10$ were rescanned one at a time,
whereas the $464{,}811{,}101$ with $n\in\{11,12\}$ were covered through the
$110{,}249{,}245$ orbit representatives of the transposed scan rather than
individually, which leaves no graph untested, since every retained encoding of a
cell $(p,q)$ is a relabelling of the $p$-side of exactly one representative and
$\alpha(G)$, the union of the maximum independent sets and $\M(G)$ are
isomorphism invariants, and the two routes were compared cell by cell wherever
both are affordable, including two cells with $n=12$. Lemma
\ref{lem:variables}, the gradedness $J(G)^t=J(G)^{(t)}$ of \cite{HHT} and Vu's
local formula \eqref{eq:local}--\eqref{eq:g} were taken as cited and not
re-derived. The independent reimplementation and its transcript accompany this
note; the program that produced Table \ref{tab:census} is retained in an
auxiliary archive available on request. The re-run is a second computation, not
a proof.

\paragraph{Scope.}
The result above is confined to bipartite graphs with at least one edge and at
most twelve nonisolated vertices. Vu's question for arbitrary bipartite graphs
remains open.

\begin{thebibliography}{9}

\bibitem{HangVu}
N.~T.~Hang and T.~Vu,
\emph{$V$-numbers of powers of cover ideals of unimodular hypergraphs},
\href{https://arxiv.org/abs/2608.18406}{arXiv:2608.18406v1} (2026).

\bibitem{HHT}
J.~Herzog, T.~Hibi, and N.~V.~Trung,
\emph{Vertex cover algebras of unimodular hypergraphs},
Proc. Amer. Math. Soc. \textbf{137} (2009), no.~2, 409--414.

\bibitem{Vu}
T.~Vu,
\emph{$V$-numbers of symbolic powers of cover ideals of graphs},
\href{https://arxiv.org/abs/2608.15268}{arXiv:2608.15268v1} (2026).

\end{thebibliography}

\end{document}